% Notes of the Montreal seminar, 1970 (cote 129, typed part): archive cote 129, site page 45 (PDF page 46).
% typescript. Draft compared: none.
% AI-transcribed by Claude Opus 5.5 (claude-opus-5-5, adaptive thinking, effort high), Message Batches API, batch mtl129_full (msgbatch_01AGnDKHwy6ATCM7ehUa6263), 2026-11-13; not yet checked by eye.
\scanpage{129/45}
\mspage{\hand{II.17}}

4.4. \uline{Cas des groupes de type multiplicatif}

On désire construire une équivalence de catégories.

\diagram{\[
\left(\begin{array}{l}
\text{p-groupes finis sur } k\\
\text{de type multiplicatif}
\end{array}\right)^{0}
\xrightarrow[\hspace{2em}]{D^{*}\ \hand{$\approx$}}
\left(\begin{array}{l}
\text{sous catégorie pleine de } \mathrm{Mod}_{\hand{$\mathcal{D}$},W\ \hand{$\ell.f.$}}\\
\text{formée des objets où } V \text{ est un}\\
\text{isomorphisme}
\end{array}\right)
\]}

La catégorie de droite est équivalente à la catégorie des
W-modules de longueur finie munis d'un isomorphisme (W-linéaire)
\[ M \xrightarrow{\ V\ } M^{\sigma} = M \otimes_{W} (W,\sigma) \]
(en effet, F est alors déterminé de façon unique par la relation FV = p).
En considérant l'isomorphisme inverse, cette catégorie est encore équiva-%
lente à la catégorie des W-modules de longueur finie munis d'un isomorphis-%
me
\[ M^{\sigma} \xrightarrow{\ V^{-1}\ } M \ . \]

Mais, d'après le corollaire 4.3.3, cette catégorie est anti-%
équivalente à la catégorie des p-groupes finis étales sur k , qui est elle-%
même anti-équivalente, par dualité de Cartier, à la catégorie des p-groupes
finis de type multiplicatif sur k .

Pour obtenir en définitive la bonne variance pour le foncteur
D* cherché, il suffit de remarquer qu'il existe une anti-équivalence de
la catégorie des p-groupes finis étales sur k avec elle-même, donnée par
la dualité de Pontryagin
\[ G \to \underline{\mathrm{Hom}}_{\mathrm{gr}}(G,\ \mathbb{Q}_{p}/\mathbb{Z}_{p}) \ . \]
\margin{x $\mathbb{Q}_p/\mathbb{Z}_p$}

\mspage{\hand{45}}
